% Nexus Scholar Course — Slide Deck Template
% Compile with: latexmk -pdf -interaction=nonstopmode <filename>.tex
%
% This template uses the metropolis Beamer theme with the custom nexus_theme.sty.
% Copy nexus_theme.sty into the same directory as this .tex file.

\documentclass[aspectratio=169]{beamer}
\usetheme{metropolis}
\usepackage{nexus_theme}

% ─── METADATA ───────────────────────────────────────────────────
\title{[Lesson Title]}
\subtitle{Module [X.Y]: [Module Name]}
\author{Nexus Scholar Suite — Modern Research Methodologies}
\date{}

\begin{document}

% ─── TITLE SLIDE ────────────────────────────────────────────────
\begin{frame}[plain]
  \titlepage
\end{frame}

% ─── LEARNING OBJECTIVES ───────────────────────────────────────
\begin{frame}{What You'll Learn}
  By the end of this lesson, you will be able to:
  \begin{enumerate}
    \item \textbf{[Bloom's verb]} [Objective 1]
    \item \textbf{[Bloom's verb]} [Objective 2]
    \item \textbf{[Bloom's verb]} [Objective 3]
  \end{enumerate}
\end{frame}

% ─── THE PROBLEM (HOOK / CONTEXT) ──────────────────────────────
\begin{frame}{The Problem}
  \begin{alertblock}{[Pain Point Title]}
    [1-3 bullet points describing the problem this lesson solves.
     Use the "old way" framing from course_review.md.]
  \end{alertblock}
  \vspace{0.3cm}
  \begin{block}{What We'll Build}
    [1-2 sentences describing the solution — the Nexus tool or technique.]
  \end{block}
\end{frame}

% ─── CONCEPT SLIDES (2-4 slides) ──────────────────────────────
\begin{frame}{[Concept 1 Title]}
  \begin{columns}[T]
    \begin{column}{0.48\textwidth}
      \begin{block}{Key Idea}
        \begin{itemize}
          \item [Point 1]
          \item [Point 2]
          \item [Point 3]
        \end{itemize}
      \end{block}
    \end{column}
    \begin{column}{0.48\textwidth}
      \begin{exampleblock}{Real-World Analogy}
        [A concrete analogy that makes the concept click.]
      \end{exampleblock}
    \end{column}
  \end{columns}
\end{frame}

% ─── ARCHITECTURE DIAGRAM (if applicable) ─────────────────────
\begin{frame}{[System / Data Flow Title]}
  \begin{center}
  \resizebox{0.85\linewidth}{!}{
  \begin{tikzpicture}[node distance=1.2cm and 1.5cm]
    \node[card] (input) {\textbf{Input}\\[description]};
    \node[primarycard, right=of input] (process) {\textbf{[Nexus Tool]}\\[key operation]};
    \node[successcard, right=of process] (output) {\textbf{Output}\\[result format]};

    \draw[flowarrow] (input) -- node[above, font=\scriptsize\bfseries] {[label]} (process);
    \draw[accentarrow] (process) -- node[above, font=\scriptsize\bfseries] {[label]} (output);
  \end{tikzpicture}
  }
  \end{center}
\end{frame}

% ─── CODE DEMONSTRATION (2-3 slides) ──────────────────────────
\begin{frame}[fragile]{[Demo Step 1: Setup]}
  \begin{lstlisting}[style=nexuspython]
# Step 1: Install and verify the toolkit
# uv run [toolkit-command] --help

from [package] import [Class]

[setup code — verified against actual source]
  \end{lstlisting}
\end{frame}

\begin{frame}[fragile]{[Demo Step 2: Core Operation]}
  \begin{lstlisting}[style=nexuspython]
# Step 2: [Description of the core operation]

[core code — verified against actual source]
  \end{lstlisting}
\end{frame}

\begin{frame}[fragile]{[Demo Step 3: Inspect Results]}
  \begin{lstlisting}[style=nexuspython]
# Step 3: Examine the output

[result inspection code]
# Expected output:
# [paste expected output here]
  \end{lstlisting}
\end{frame}

% ─── KEY TAKEAWAY ─────────────────────────────────────────────
\begin{frame}{Key Takeaway}
  \begin{block}{What We Accomplished}
    \begin{enumerate}
      \item [Takeaway 1 — maps to Objective 1]
      \item [Takeaway 2 — maps to Objective 2]
      \item [Takeaway 3 — maps to Objective 3]
    \end{enumerate}
  \end{block}
  \vspace{0.3cm}
  \textbf{Next:} Module [X.Y] — [Next Lesson Title]
\end{frame}

% ─── EXERCISE PREVIEW ─────────────────────────────────────────
\begin{frame}{Your Turn: Exercise}
  \begin{exampleblock}{Exercise: [Title]}
    [1-2 sentence description of what the student will do in the exercise.]
  \end{exampleblock}
\end{frame}

% ─── STANDOUT CLOSE ───────────────────────────────────────────
\begin{frame}[standout]
  [A single memorable sentence that captures the lesson's essence.]
\end{frame}

\end{document}
